\section{Numerical Linear Algebra (NLA) \& Direct Solvers}

Solving linear systems $A x = b$ is the fundamental computational engine of scientific computing. We study how matrix factorizations transform dense, intractable equations into easily solvable triangular forms while preserving numerical stability.

% -----------------------------------------------------------------------------
\subsection{Conditioning of Linear Systems \& The Golden Rule}

\begin{definition}
    For an invertible matrix $A \in \mathbb{R}^{n \times n}$, the \textbf{matrix condition number} with respect to a matrix norm is:
    \begin{equation}
        \operatorname{cond}(A) := \|A\| \cdot \|A^{-1}\| \ge 1.
    \end{equation}
    In the Euclidean 2-norm, this is the ratio of extreme singular values:
    \begin{equation}
        \operatorname{cond}_2(A) = \frac{\sigma_{\max}(A)}{\sigma_{\min}(A)} = \frac{\sigma_1}{\sigma_n}.
    \end{equation}
\end{definition}

For a perturbed system $(A + \Delta A)(x + \Delta x) = b + \Delta b$, the relative error in $x$ satisfies:
\begin{equation}
    \frac{\|\Delta x\|}{\|x\|} \le \frac{\operatorname{cond}(A)}{1 - \operatorname{cond}(A) \frac{\|\Delta A\|}{\|A\|}} \left( \frac{\|\Delta A\|}{\|A\|} + \frac{\|\Delta b\|}{\|b\|} \right).
\end{equation}
Solving $Ax = b$ in IEEE double precision loses roughly $\log_{10}(\operatorname{cond}(A))$ decimal digits. If $\operatorname{cond}(A) \ge 10^{16}$, all precision is lost.

\begin{important}
    \textbf{Never compute the explicit matrix inverse $A^{-1}$ to solve $Ax = b$!} Computing $A^{-1}$ requires $2n^3$ flops (three times more operations than LU factorization), destroys backward numerical stability if $A$ is ill-conditioned, and turns sparse matrices into completely dense matrices ($\mathcal{O}(n^2)$ memory). Always solve $Ax = b$ via triangular factorizations.
\end{important}

% -----------------------------------------------------------------------------
\subsection{LU \& Cholesky Factorizations}

In \textbf{LU factorization with partial pivoting}, any invertible matrix $A \in \mathbb{R}^{n \times n}$ can be factored as:
\begin{equation}
    P A = L U,
\end{equation}
where $P$ is a permutation matrix, $L$ is unit lower-triangular ($L_{i,i} = 1, |L_{i,j}| \le 1$), and $U$ is upper-triangular. Partial pivoting searches column $k$ for the entry of maximum magnitude and swaps rows, ensuring $|L_{i,j}| \le 1$ to prevent exponential roundoff growth. The factorization takes $\sum_{k=1}^{n-1} 2(n-k)^2 \approx \frac{2}{3}n^3$ flops. Solving proceeds via forward substitution $L y = P b$ followed by backward substitution $U x = y$ in $2n^2$ flops. For multiple right-hand sides, $A = LU$ is computed once, and each subsequent solve costs only $2n^2$ flops.

A matrix $A \in \mathbb{R}^{n \times n}$ is \textbf{Symmetric Positive Definite (SPD)} if $A = A^T$ and $x^T A x > 0$ for all $x \neq 0$. The \textbf{Cholesky factorization} decomposes every SPD matrix uniquely into:
\begin{equation}
    A = L L^T \quad (\text{or } A = R^T R \text{ with upper-triangular } R = L^T),
\end{equation}
where $L$ is lower-triangular with strictly positive diagonal entries $L_{i,i} > 0$. Because $A_{i,i} = \sum_{k=1}^i L_{i,k}^2$, the entries of $L$ satisfy $|L_{i,k}| \le \sqrt{A_{i,i}}$, guaranteeing unconditional numerical stability \textbf{without pivoting}. Cholesky requires only \textbf{$\frac{1}{3}n^3$ flops} (half the cost of LU) and stores only half of the matrix entries.

% -----------------------------------------------------------------------------
\subsection{QR Factorization \& Orthogonal Transformations}

A square matrix $Q \in \mathbb{R}^{m \times m}$ is orthogonal if $Q^T Q = Q Q^T = I$. Orthogonal transformations preserve Euclidean norms ($\|Q x\|_2 = \|x\|_2$) and condition numbers ($\operatorname{cond}_2(Q A) = \operatorname{cond}_2(A)$).

Any rectangular matrix $A \in \mathbb{R}^{m \times n}$ ($m \ge n$) admits an \textbf{economic (reduced) QR factorization}:
\begin{equation}
    A = Q_1 R_1,
\end{equation}
where $Q_1 \in \mathbb{R}^{m \times n}$ has orthonormal columns ($Q_1^T Q_1 = I_n$) and $R_1 \in \mathbb{R}^{n \times n}$ is upper-triangular. For dense matrices, \textbf{Householder reflections} $H = I - 2 \frac{v v^T}{v^T v}$ reflect a vector across a hyperplane onto $\|x\|_2 e_1$, zeroing out all subdiagonal entries in a column simultaneously in $2m n^2 - \frac{2}{3}n^3$ flops while preserving machine-precision orthogonality. Alternatively, \textbf{Givens rotations} zero out individual matrix entries selectively, making them ideal for sparse or banded matrices.

% -----------------------------------------------------------------------------
\subsection{Singular Value Decomposition (SVD)}

\begin{theorem}
    In the \textbf{Singular Value Decomposition (SVD)}, every matrix $A \in \mathbb{R}^{m \times n}$ can be factored as:
    \begin{equation}
        A = U \Sigma V^T = \sum_{j=1}^{p} \sigma_j u_j v_j^T, \quad p = \min(m, n),
    \end{equation}
    where $U \in \mathbb{R}^{m \times m}$ and $V \in \mathbb{R}^{n \times n}$ are orthogonal, and $\Sigma$ contains singular values $\sigma_1 \ge \sigma_2 \ge \dots \ge \sigma_p \ge 0$.
\end{theorem}

Geometrically, $A$ maps the unit sphere in $\mathbb{R}^n$ into a hyper-ellipsoid in $\mathbb{R}^m$. The right singular vectors $v_j$ map to the principal semi-axes, whose lengths are the singular values $\sigma_j$, and whose directions are given by the left singular vectors $u_j$.

\begin{theorem}
    By the \textbf{Eckart-Young-Mirsky theorem}, for any $1 \le k \le \operatorname{rank}(A)$, the truncated SVD sum $A_k := \sum_{j=1}^{k} \sigma_j u_j v_j^T$ is the unique best rank-$k$ approximation of $A$ under both the spectral and Frobenius norms:
    \begin{equation}
        \min_{\operatorname{rank}(B) \le k} \|A - B\|_2 = \|A - A_k\|_2 = \sigma_{k+1}, \qquad \min_{\operatorname{rank}(B) \le k} \|A - B\|_F = \|A - A_k\|_F = \sqrt{\sum_{j=k+1}^{p} \sigma_j^2}.
    \end{equation}
\end{theorem}

The numerical rank of a matrix is the count of singular values satisfying $\sigma_j > \varepsilon_{\text{mach}} \sigma_1$. SVD is the only unconditionally robust tool for numerical rank determination.

% -----------------------------------------------------------------------------
\subsection{Sparse Matrix Representations}

When $A \in \mathbb{R}^{n \times n}$ has only $\mathcal{O}(n)$ non-zero entries, storing all $n^2$ entries would exhaust memory. Sparse formats store only non-zero values. The \textbf{Coordinate Format (COO)} stores arrays \texttt{(row, col, data)} with simple $\mathcal{O}(1)$ element insertion, used primarily for matrix assembly. The \textbf{Compressed Sparse Row (CSR)} format stores \texttt{data}, column \texttt{indices}, and row pointers \texttt{indptr}, optimized for matrix-vector products $A x$ in $\mathcal{O}(\operatorname{nnz})$ flops and row slicing. The \textbf{Compressed Sparse Column (CSC)} format is the column-oriented counterpart of CSR, optimized for column slicing and used internally by direct sparse solvers.

During LU or Cholesky elimination, zero entries in $A$ frequently become non-zero in $L$ and $U$, a phenomenon called \textbf{fill-in}. Permuting rows and columns ($P A P^T$) prior to factorization mitigates fill-in. Heuristics like Approximate Minimum Degree (AMD) eliminate low-degree vertices first, while Reverse Cuthill-McKee (RCM) reduces matrix bandwidth.
